\section{Introduction}

Hundreds of millions of South Asian language speakers communicate with digital
systems---including large language models (LLMs)---not in their native script,
but in romanized Latin characters. A Hindi speaker types \textit{aap kaise hain}
rather than \textit{आप कैसे हैं}; a Telugu speaker types \textit{meeru ela
unnaru} rather than \textit{మీరు ఎలా ఉన్నారు}. This practice is not a niche
behavior. It is the dominant input mode on smartphones, where QWERTY keyboards
are universal and native-script input remains cumbersome. Prevalence studies in
real-world messaging corpora place romanized usage at 83\% for Hindi and 73\%
for Telugu~\cite{khullar2025scriptgap}; for hundreds of millions of users,
romanized Latin script \emph{is} the de facto interface for their native
language.

Despite its scale, this modality is nearly absent from standard NLP evaluation.
Benchmarks such as XQuAD, IndicQA, and MLQA use formal native script.
Language detectors routinely misclassify romanized Hindi or Telugu as English,
preventing correct labeling. Tokenizers trained on native script have no learned
representations for romanized equivalents: a word like \textit{untundi} is split
into arbitrary English subword units with no semantic grounding. Datasets such as
Dakshina~\cite{roark2020dakshina} and Aksharantar~\cite{madhani2023aksharantar}
document the scope of Latin-script South Asian language processing and the effort
needed for robust transliteration, but large-scale romanized \emph{generation}
benchmarks---especially for open-ended, long-form responses---remain absent.

This invisibility has a subtler, more consequential dimension: automatic
evaluation metrics may fail to detect quality failures that native speakers
readily identify. A model response that mixes native-script characters
mid-sentence, uses an awkward register, or chooses semantically wrong words can
still receive a high BERTScore if the multilingual encoder has no meaningful
representation of romanized text. We show that even BGE-M3 cosine
similarity---a state-of-the-art cross-lingual embedding model specifically
designed for multilingual retrieval~\cite{chen2024bgem3}---produces scores in
the 0.76--0.84 range across all model and language combinations, masking
systematic quality differences that native speakers consistently detect. Prior
work has shown that romanized input degrades LLM performance on structured tasks
by up to 24 points relative to native script~\cite{khullar2025scriptgap}; our
finding is that this degradation is further compounded by evaluation frameworks
that cannot see it.

We address this gap with a controlled multi-model evaluation. Using 142 English
query--response pairs from the OpenAssistant dataset, we prompt four
LLMs---Claude (Anthropic), GPT-4o (OpenAI), Gemini (Google), and Qwen3.5-35B
served locally via Ollama---to produce romanized Hindi, Telugu, and Kannada
versions. We evaluate outputs with multilingual BERTScore~\cite{zhang2020bertscore}
and BGE-M3 cosine similarity under a failure-aware scoring policy that assigns
zero to empty, copied, and native-script-corrupted outputs rather than excluding
them, ensuring that reliability and semantic quality are reported together. We
include Korean as a negative control: Korean speakers do not romanize in everyday
communication, providing a baseline that isolates romanization---not non-English
use alone---as the source of any observed degradation. A native-speaker human
evaluation on sampled responses per language completes the assessment, directly
measuring the gap between automatic scores and human judgment.

Our contributions are:
\begin{itemize}
    \item a multi-model, multi-language romanized generation benchmark covering
    three Indic languages (Hindi, Telugu, Kannada) across four LLMs, with Korean
    as a typological negative control;

    \item a failure-aware scoring framework that treats empty, copied, and
    script-corrupted outputs as failures rather than omissions, giving a
    complete picture of model reliability;

    \item empirical evidence that state-of-the-art automatic metrics---including
    BGE-M3 embedding similarity---cannot detect quality failures that native
    speakers consistently identify, quantifying the evaluation blind spot; and

    \item a native-speaker human evaluation rubric and annotated sample designed
    to support future benchmarking of romanized Indic generation.
\end{itemize}
